\documentclass[12pt,a4paper]{article}

\usepackage[T2A]{fontenc}
\usepackage[utf8]{inputenc}
\usepackage[russian]{babel}
\usepackage{geometry}
\usepackage{hyperref}
\usepackage{booktabs}
\usepackage{listings}
\usepackage{xcolor}

\geometry{margin=2cm}
\hypersetup{colorlinks=true, linkcolor=blue, urlcolor=blue, citecolor=blue}

\lstset{
  basicstyle=\ttfamily\small,
  breaklines=true,
  columns=fullflexible,
  frame=single,
  rulecolor=\color{black!20},
}

\begin{document}

\begin{titlepage}
  \begin{center}
    {\large Университет ИТМО}\\[4mm]
    {\large Факультет/направление: \underline{\hspace{7cm}}}\\[10mm]

    {\LARGE \textbf{Лабораторная работа №2}}\\[2mm]
    {\Large \textbf{Блочный кэш страниц (VT Page Cache)}}\\[6mm]

    {\large Вариант: \textbf{Optimal (с подсказкой)}}\\[12mm]
  \end{center}

  \vfill

  \begin{flushleft}
    Студент: \textbf{Брагин Роман}, группа \textbf{P3316}\\
    Преподаватель: \textbf{Сорокин Артём Николаевич}\\
  \end{flushleft}

  \vfill
  \begin{center}
    Санкт-Петербург, \the\year
  \end{center}
\end{titlepage}

\section{Текст задания (кратко)}
Необходимо реализовать блочный кэш страниц в user-space в виде динамической библиотеки на C.
Библиотека должна предоставлять API, аналогичный системным вызовам: открытие/закрытие файла, чтение, запись, позиционирование, синхронизация.
Операции чтения/записи на диск должны выполняться в обход page cache ОС.

Вариант: \textbf{Optimal}. Для него требуется предоставить пользователю возможность подсказать момент следующего обращения к данным (например, через отдельную функцию \texttt{advice}).

\section{Обзор реализации}
\subsection{Структура проекта}
\begin{itemize}
  \item \texttt{lib/vtpc.c}, \texttt{lib/vtpc.h} --- реализация библиотеки кэша страниц.
  \item \texttt{test/} --- тесты корректности (поведение сравнивается с \texttt{libc}).
  \item \texttt{bench/io-load.c} --- IO-нагрузчик для сравнения производительности \texttt{libc} и \texttt{vtpc}.
\end{itemize}

\subsection{Реализованный API}
Библиотека предоставляет функции:
\begin{itemize}
  \item \texttt{vtpc\_open}, \texttt{vtpc\_close}, \texttt{vtpc\_read}, \texttt{vtpc\_write}, \texttt{vtpc\_lseek}, \texttt{vtpc\_fsync}
  \item \texttt{vtpc\_advice(fd, offset, next\_use)} --- подсказка для Optimal
  \item алиасы \texttt{lab2\_*} для соответствия формату задания (например, \texttt{lab2\_open}, \texttt{lab2\_advice})
\end{itemize}

\subsection{Кэш страниц}
\begin{itemize}
  \item Кэш состоит из фиксированного числа страниц (\texttt{VTPC\_CACHE\_PAGES}), каждая страница имеет размер (\texttt{VTPC\_PAGE\_SIZE}).
  \item Данные страниц хранятся в выровненных буферах (через \texttt{posix\_memalign}) для корректной работы с \texttt{O\_DIRECT}.
  \item При \texttt{read/write} обрабатываются произвольные размеры и смещения: операция раскладывается на затронутые страницы, а работа с диском выполняется страничными блоками.
  \item Модифицированные страницы помечаются как \texttt{dirty} и сбрасываются на диск при вытеснении/ \texttt{fsync}/ \texttt{close}.
\end{itemize}

\subsection{Обход page cache ОС}
При открытии файла библиотека пытается использовать \texttt{O\_DIRECT}.
Если файловая система не поддерживает \texttt{O\_DIRECT}, библиотека автоматически открывает файл без него, при этом старается минимизировать влияние page cache ОС (best-effort через \texttt{posix\_fadvise(..., DONTNEED)}).

\subsection{Политика вытеснения Optimal}
\begin{itemize}
  \item Для каждой страницы поддерживается значение \texttt{next\_use} --- ``время'' следующего обращения.
  \item При необходимости вытеснения выбирается страница с максимальным \texttt{next\_use} (Belady/Optimal).
  \item При равенстве \texttt{next\_use} используется тай-брейк по давности доступа (LRU).
  \item Пользователь задаёт \texttt{next\_use} через \texttt{vtpc\_advice/lab2\_advice}.
\end{itemize}

\section{Проверка корректности}
Корректность проверялась тестами из каталога \texttt{test/}:
\begin{itemize}
  \item \texttt{test\_basic}
  \item \texttt{test\_seq}
  \item \texttt{test\_random}
\end{itemize}

Тесты запускаются через CTest:
\begin{lstlisting}
make test
# или
ctest --test-dir build/test --output-on-failure
\end{lstlisting}

\section{Замеры производительности (до/после)}
\subsection{Методика}
Для замеров использовался нагрузчик \texttt{bench/io-load} со следующими параметрами:
\begin{itemize}
  \item блок: 4096 байт
  \item количество операций: 20000
  \item диапазон: 0--262144 (256 КиБ, горячий набор данных)
  \item порядок: random
  \item сравнение:
    \begin{itemize}
      \item \texttt{libc}: \texttt{pread/pwrite} с \texttt{O\_DIRECT=on}
      \item \texttt{vtpc}: библиотека кэша, \texttt{VTPC\_PAGE\_SIZE=4096}, \texttt{VTPC\_CACHE\_PAGES=64}
    \end{itemize}
\end{itemize}

\subsection{Результаты}
\begin{table}[h!]
  \centering
  \begin{tabular}{@{}lrrr@{}}
    \toprule
    Сценарий & libc O\_DIRECT (MB/s) & vtpc (MB/s) & Ускорение \\
    \midrule
    Random read (256 KiB hot set)  & 21.60  & 5032.23 & \(\approx 233\times\) \\
    Random write (256 KiB hot set) & 17.51  & 4464.67 & \(\approx 255\times\) \\
    \bottomrule
  \end{tabular}
  \caption{Пример сравнения производительности до/после внедрения page cache}
\end{table}

\textbf{Комментарий.} На ``горячем'' наборе данных (256 КиБ) кэш позволяет практически полностью убрать повторные обращения к диску: страницы читаются/записываются в RAM, а на диск сбрасываются агрегированно.

\section{Вывод}
В рамках работы реализован пользовательский page cache в виде динамической библиотеки на C с обходом page cache ОС и политикой вытеснения Optimal (Belady) с подсказками следующего обращения.
Тесты корректности проходят, а на сценариях с повторным доступом к одному и тому же диапазону данных наблюдается существенный выигрыш по производительности.

\end{document}


